% GENERATED FILE. Edit canonical lessons, then run npm run generate:books.
% canonical-source-hash: fnv1a64:8e6c713a0e396b32
% canonical-lessons: JA-C434-kansha-suru, JA-C434-doujou-suru, JA-C434-kinodoku, JA-C434-mijime, JA-C434-hagemasu

\chapter{Thanks and Sympathy}
\label{ch:ja-thanks-and-sympathy}
\hlchaptermodality{434}

\begin{chapteropening}
\textbf{By the end of this chapter:} \emph{I can thank someone and show sympathy.}

Complete the last lesson of chapter 434: I can thank someone and show sympathy.
\end{chapteropening}

\section[\texorpdfstring{\ja{かんしゃする} (kansha suru)}{かんしゃする (kansha suru)}]{\ja{かんしゃする} (kansha suru) — to be grateful; to thank}

\label{lesson:JA-C434-kansha-suru}



\begin{quote}
Before the new one: say the Japanese for an interest, then the Japanese for admiration.
\end{quote}

\subsection*{You'll want to know: \ja{かんしゃする}}
\textbf{\ja{かんしゃする}} — \emph{kansha suru} — \textquotedblleft{}to be grateful; to thank\textquotedblright{}.

\begin{quote}
\textbf{\ja{りょうしんに} \ja{かんしゃしています。}} — \emph{ryōshin ni kansha shite imasu} — \textquotedblleft{}I am grateful to my parents.\textquotedblright{}
\end{quote}

Polite \textbf{\ja{かんしゃします}}, \emph{kanshashimasu}; a \ja{する} verb. Past (polite): \textbf{\ja{がくせいが} \ja{せんせいに} \ja{かんしゃしました。}}, \emph{gakusei ga sensei ni kansha shimashita}.

\begin{cousinweb}
\textbf{\ja{かんしゃする}} (\emph{kansha suru}) is Sino-Japanese (kango) \emph{kansha} with the native verb \emph{suru}, \textquotedblleft{}to do\textquotedblright{}: two Chinese characters read with their on'yomi, readings borrowed from Chinese. They are \emph{kan}, \textquotedblleft{}emotion\textquotedblright{} (Middle Chinese \emph{komX}) and \emph{sha}, \textquotedblleft{}thank\textquotedblright{} (\emph{zjaeH}) (KANJIDIC2; Middle Chinese in Baxter's transcription).
\end{cousinweb}

\begin{grammarlens}[title={What you've built}]
One more word for describing what you lived and felt.
\end{grammarlens}

\subsection*{Your turn}
\begin{itemize}
\raggedright
  \item \emph{Say it:} \emph{kansha suru}
  \item \emph{Say it:} \emph{kansha suru}, once more
  \item \emph{Say it:} \emph{ryōshin ni kansha shite imasu} — I am grateful to my parents.
  \item \emph{Recall:} say the Japanese for an interest, then the Japanese for admiration, then say \emph{kansha suru} again
\end{itemize}

\begin{tcolorbox}[breakable,colback=teal!4,colframe=teal!35!black,title={Before you move on}]
What does \ja{かんしゃする} mean? (\textquotedblleft{}To be grateful; to thank\textquotedblright{}.) Say it once more.
\end{tcolorbox}
\section[\texorpdfstring{\ja{どうじょうする} (dōjō suru)}{どうじょうする (dōjō suru)}]{\ja{どうじょうする} (dōjō suru) — to feel sympathy}

\label{lesson:JA-C434-doujou-suru}



\begin{quote}
Before the new one: say the Japanese for admiration, then the Japanese for to be grateful.
\end{quote}

\subsection*{You'll want to know: \ja{どうじょうする}}
\textbf{\ja{どうじょうする}} — \emph{dōjō suru} — \textquotedblleft{}to feel sympathy\textquotedblright{}.

\begin{quote}
\textbf{\ja{かれが} \ja{けがを} \ja{して、みんなが} \ja{どうじょうしました。}} — \emph{kare ga kega o shite, minna ga dōjō shimashita} — \textquotedblleft{}He got hurt, and everyone felt for him.\textquotedblright{}
\end{quote}

Polite \textbf{\ja{どうじょうします}}, \emph{dōjōshimasu}; a \ja{する} verb. Past (polite): \textbf{\ja{ともだちが} \ja{どうじょうしました。}}, \emph{tomodachi ga dōjō shimashita}.

\begin{cousinweb}
\textbf{\ja{どうじょうする}} (\emph{dōjō suru}) is Sino-Japanese (kango) \emph{dōjō} with the native verb \emph{suru}, \textquotedblleft{}to do\textquotedblright{}: two Chinese characters read with their on'yomi, readings borrowed from Chinese. They are \emph{dō}, \textquotedblleft{}same\textquotedblright{} (Middle Chinese \emph{duwng}) and \emph{jō}, \textquotedblleft{}sympathy\textquotedblright{} (\emph{dzjeng}) (KANJIDIC2; Middle Chinese in Baxter's transcription).
\end{cousinweb}

\begin{grammarlens}[title={What you've built}]
One more word for describing what you lived and felt.
\end{grammarlens}

\subsection*{Your turn}
\begin{itemize}
\raggedright
  \item \emph{Say it:} \emph{dōjō suru}
  \item \emph{Say it:} \emph{dōjō suru}, once more
  \item \emph{Say it:} \emph{kare ga kega o shite, minna ga dōjō shimashita} — He got hurt, and everyone felt for him.
  \item \emph{Recall:} say the Japanese for admiration, then the Japanese for to be grateful, then say \emph{dōjō suru} again
\end{itemize}

\begin{tcolorbox}[breakable,colback=teal!4,colframe=teal!35!black,title={Before you move on}]
What does \ja{どうじょうする} mean? (\textquotedblleft{}To feel sympathy\textquotedblright{}.) Say it once more.
\end{tcolorbox}
\section[\texorpdfstring{\ja{きのどく} (kinodoku)}{きのどく (kinodoku)}]{\ja{きのどく} (kinodoku) — pitiable; unfortunate}

\label{lesson:JA-C434-kinodoku}



\begin{quote}
Before the new one: say the Japanese for to be grateful, then the Japanese for to feel sympathy.
\end{quote}

\subsection*{You'll want to know: \ja{きのどく}}
\textbf{\ja{きのどく}} — \emph{kinodoku} — \textquotedblleft{}pitiable; unfortunate\textquotedblright{}.

\begin{quote}
\textbf{\ja{ともだちは} \ja{さいふを} \ja{なくして、きのどくでした。}} — \emph{tomodachi wa saifu o nakushite, kinodoku deshita} — \textquotedblleft{}My friend lost his wallet, poor thing.\textquotedblright{}
\end{quote}

A na-adjective: \ja{な} before a noun.

\begin{cousinweb}
\textbf{\ja{きのどく}} (\emph{kinodoku}) is a mixed phrase: Sino-Japanese \emph{ki}, \textquotedblleft{}spirit\textquotedblright{} (Middle Chinese \emph{khj+jH}), and \emph{doku}, \textquotedblleft{}poison\textquotedblright{} (\emph{dowk}), both borrowed from Chinese, joined by the native particle \emph{no}, \textquotedblleft{}of\textquotedblright{} (JMdict, KANJIDIC2; Middle Chinese in Baxter's transcription).
\end{cousinweb}

\begin{grammarlens}[title={What you've built}]
One more word for describing what you lived and felt.
\end{grammarlens}

\subsection*{Your turn}
\begin{itemize}
\raggedright
  \item \emph{Say it:} \emph{kinodoku}
  \item \emph{Say it:} \emph{kinodoku}, once more
  \item \emph{Say it:} \emph{tomodachi wa saifu o nakushite, kinodoku deshita} — My friend lost his wallet, poor thing.
  \item \emph{Recall:} say the Japanese for to be grateful, then the Japanese for to feel sympathy, then say \emph{kinodoku} again
\end{itemize}

\begin{tcolorbox}[breakable,colback=teal!4,colframe=teal!35!black,title={Before you move on}]
What does \ja{きのどく} mean? (\textquotedblleft{}Pitiable; unfortunate\textquotedblright{}.) Say it once more.
\end{tcolorbox}
\section[\texorpdfstring{\ja{みじめ} (mijime)}{みじめ (mijime)}]{\ja{みじめ} (mijime) — miserable; wretched}

\label{lesson:JA-C434-mijime}



\begin{quote}
Before the new one: say the Japanese for to feel sympathy, then the Japanese for pitiable.
\end{quote}

\subsection*{You'll want to know: \ja{みじめ}}
\textbf{\ja{みじめ}} — \emph{mijime} — \textquotedblleft{}miserable; wretched\textquotedblright{}.

\begin{quote}
\textbf{\ja{あめの} \ja{なかの} \ja{ねこは} \ja{みじめでした。}} — \emph{ame no naka no neko wa mijime deshita} — \textquotedblleft{}The cat out in the rain was miserable.\textquotedblright{}
\end{quote}

A na-adjective: \ja{な} before a noun.

\begin{cousinweb}
\textbf{\ja{みじめ}} (\emph{mijime}) is a native Japanese word (wago), inherited rather than borrowed from Chinese; its kanji spelling was attached to it later (JMdict).
\end{cousinweb}

\begin{grammarlens}[title={What you've built}]
One more word for describing what you lived and felt.
\end{grammarlens}

\subsection*{Your turn}
\begin{itemize}
\raggedright
  \item \emph{Say it:} \emph{mijime}
  \item \emph{Say it:} \emph{mijime}, once more
  \item \emph{Say it:} \emph{ame no naka no neko wa mijime deshita} — The cat out in the rain was miserable.
  \item \emph{Recall:} say the Japanese for to feel sympathy, then the Japanese for pitiable, then say \emph{mijime} again
\end{itemize}

\begin{tcolorbox}[breakable,colback=teal!4,colframe=teal!35!black,title={Before you move on}]
What does \ja{みじめ} mean? (\textquotedblleft{}Miserable; wretched\textquotedblright{}.) Say it once more.
\end{tcolorbox}
\section[\texorpdfstring{\ja{はげます} (hagemasu)}{はげます (hagemasu)}]{\ja{はげます} (hagemasu) — to encourage; to cheer up}

\label{lesson:JA-C434-hagemasu}



\begin{quote}
Before the new one: say the Japanese for pitiable, then the Japanese for miserable.
\end{quote}

\subsection*{You'll want to know: \ja{はげます}}
\textbf{\ja{はげます}} — \emph{hagemasu} — \textquotedblleft{}to encourage; to cheer up\textquotedblright{}.

\begin{quote}
\textbf{\ja{ともだちが} \ja{おちこんだので、はげましました。}} — \emph{tomodachi ga ochikonda node, hagemashimashita} — \textquotedblleft{}My friend was feeling down, so I cheered him up.\textquotedblright{}
\end{quote}

Polite \textbf{\ja{はげまします}}, \emph{hagemashimasu}; a godan verb. Past (polite): \textbf{\ja{ははが} \ja{むすこを} \ja{はげましました。}}, \emph{haha ga musuko o hagemashimashita}.

\begin{cousinweb}
\textbf{\ja{はげます}} (\emph{hagemasu}) is a native Japanese word (wago), inherited rather than borrowed from Chinese; its kanji spelling was attached to it later (JMdict).
\end{cousinweb}

\begin{grammarlens}[title={What you've built}]
One more word for describing what you lived and felt.
\end{grammarlens}

\subsection*{Your turn}
\begin{itemize}
\raggedright
  \item \emph{Say it:} \emph{hagemasu}
  \item \emph{Say it:} \emph{hagemasu}, once more
  \item \emph{Say it:} \emph{tomodachi ga ochikonda node, hagemashimashita} — My friend was feeling down, so I cheered him up.
  \item \emph{Recall:} say the Japanese for pitiable, then the Japanese for miserable, then say \emph{hagemasu} again
\end{itemize}

\begin{tcolorbox}[breakable,colback=teal!4,colframe=teal!35!black,title={Before you move on}]
What does \ja{はげます} mean? (\textquotedblleft{}To encourage; to cheer up\textquotedblright{}.) Say it once more.
\end{tcolorbox}
